\documentclass[11pt]{article}
\usepackage[margin=1in]{geometry}
\usepackage{amsmath,amssymb,amsthm,booktabs,longtable,array}
\usepackage{microtype,listings,xcolor}
\setlength{\emergencystretch}{1em}
\usepackage[hidelinks]{hyperref}
\usepackage{url}
\definecolor{ink}{HTML}{24384B}
\lstdefinestyle{reference}{language=Python,basicstyle=\ttfamily\footnotesize,
 keywordstyle=\bfseries\color{ink},commentstyle=\itshape\color{ink},
 columns=fullflexible,keepspaces=true,showstringspaces=false,
 breaklines=true,breakatwhitespace=true,frame=lines,
 aboveskip=0.7\baselineskip,belowskip=0.7\baselineskip}
\newcommand{\ins}[1]{\texttt{#1}}
\newcommand{\T}[1]{\operatorname{T}_{#1}}
\newcommand{\N}[1]{\operatorname{RN}_{#1}}
\newcommand{\R}{\operatorname{R}_{\mathrm{RC}}}
\newcommand{\hex}[1]{\texttt{#1}}
\newtheorem{proposition}{Proposition}
\hypersetup{pdftitle={Reconstructing FSIN, FCOS and FSINCOS from Public Constants and Processor Tests},
 pdfauthor={},pdfsubject={Executable behavioral reconstruction and validation}}
\input{generated-suite/numbers.tex}
\title{Reconstructing FSIN, FCOS and FSINCOS\\
from Public Constants and Processor Tests}
\author{}
\date{September 2026}

\begin{document}
\maketitle
\begin{abstract}
The constants used to approximate sine and cosine do not, by themselves,
explain the results of x87 instructions. The order of operations and the
rounding of intermediate values also matter. We reconstruct executable
algorithms for \ins{FSIN}, \ins{FCOS} and \ins{FSINCOS} on the tested Intel
Skylake Core i7 and Xeon processors, starting from Ken Shirriff's published
Pentium ROM analysis. The main findings are distinct polynomial schedules
for standalone and paired instructions, unequal operand widths in the
standalone multiplication rule, and different rounding destinations within
the shared table calculation. Exact arithmetic and measurements near rounding boundaries constrain these
details through the final result and C1 bit. The resulting programs use fixed arithmetic rules and the published
constants, and match the reported hardware tests and saved-data replays.
Bounds on reduction and intermediate widths support their implementation;
they do not establish a unique physical circuit or agreement for every input.
The accompanying reference also covers \ins{FPTAN}, \ins{FPATAN},
\ins{F2XM1}, \ins{FYL2X} and \ins{FYL2XP1}, with complete pseudocode,
constants and programmer interfaces for all eight instructions.
\end{abstract}

\section{Introduction}

An emulator that replaces \ins{FSIN} with a library sine function can return
different bits from the processor. The library's answer may be closer to the
mathematical sine, but it can still be wrong for a program that expects x87
behavior. Intel describes this distinction and the limits of the instructions'
range reduction~\cite{intelDifference}. Our aim is to reproduce the processor's
calculation, including the rounding decisions that affect its last bits.

The main contribution of this work is the reconstruction of \ins{FSIN},
\ins{FCOS} and \ins{FSINCOS}. Published constants explain the approximation
families. Reproducing the instruction results also requires the order of
operations, the precision of each intermediate value, and the final rounding
rule. Standalone sine and cosine and the paired instruction use different
finite-precision evaluations of the same approximation data.

We specify those calculations as explicit arithmetic programs. Standalone
\ins{FSIN}/\ins{FCOS} split the polynomial into odd and even coefficient
chains; \ins{FSINCOS} uses two Horner chains sharing a square. The standalone
multiply retains different widths for its two operands. The shared table
calculation also contains a product that rounds directly to 64 bits, while
ordinary products truncate to 67. The final additions, sign restoration and
C1 rule are specified as part of the algorithm. Together, these rules explain how the approximation data produce the
observed result bits and C1.

The evidence combines inputs selected to distinguish arithmetic rules,
independent exact implementations, completed native runs and extensive replays
of earlier hardware observations. We also establish properties of the
reconstructed programs, including exact reduction with the finite divisor and
bounds on intermediate values. These serve different purposes: hardware
measurements select numerical behavior, while the bounds check that the code
can carry out the specified calculation without losing additional bits.

The paper is organized into three families. Section~\ref{family:trig}
contains the main reconstruction: \ins{FSIN}, \ins{FCOS} and \ins{FSINCOS},
from their mathematical basis through the reconstructed arithmetic and its
validation. Section~\ref{family:tangent} covers \ins{FPTAN} and \ins{FPATAN},
which followed comparatively directly from the published material.
Section~\ref{family:explog} places the shorter exponential and logarithm
implementations last. Each family has its own algorithm, constants and
evidence discussion; the appendices follow the same order.

\begin{table}[ht]
\centering\small
\begin{tabular}{@{}lll@{}}
\toprule
Instruction & Mathematical function & Family\\
\midrule
\ins{FSIN} & $\sin x$ & Sine and cosine\\
\ins{FCOS} & $\cos x$ & Sine and cosine\\
\ins{FSINCOS} & $(\sin x,\cos x)$ & Sine and cosine\\
\midrule
\ins{FPTAN} & $\tan x$, with a pushed result & Tangent and arctangent\\
\ins{FPATAN} & $\operatorname{atan2}(y,x)$ & Tangent and arctangent\\
\midrule
\ins{F2XM1} & $2^x-1$ & Exponential and logarithms\\
\ins{FYL2X} & $y\log_2x$ & Exponential and logarithms\\
\ins{FYL2XP1} & $y\log_2(1+x)$ & Exponential and logarithms\\
\bottomrule
\end{tabular}
\caption{The main reconstruction and the accompanying implementations.
All eight instructions have complete specifications~\cite{intelSdm}.
Their mathematical functions describe what they approximate; the models
also preserve the processor's intermediate rounding.}
\label{tab:scope}
\end{table}

\section{FSIN, FCOS and FSINCOS}
\label{family:trig}

The sine and cosine reconstruction starts with the published approximation
methods and determines the finite-precision operations used to evaluate them.

\subsection{Published constants and implementations}

Shirriff decoded the Pentium floating-point ROM and explained how its
constants are used~\cite{shirriff2025rom}. That work supplies the main constants
and many of the algorithmic ideas used here. We build on it by working out the
precision and rounding steps needed to reproduce the tested processors.
Appendix~\ref{app:trig-constants} records the literal values, including two
cosine-table words that differ from the saved published transcription.

Harrison formally verified trigonometric algorithms used for IA-32
compatibility on Itanium~\cite{harrison2000}. Intel's assembly implementation
is preserved in glibc 2.38~\cite{glibc238}; our repository includes a separate
translation that uses integer arithmetic. The IA-64 literature also explains
the underlying reduction and approximation methods~\cite{harrison1999}.
These sources help explain the algorithms, but we do not assume that their
rounding steps match Skylake.

\subsection{Measurements and exact arithmetic}

Each hardware observation records the instruction, raw input bits, rounding
control (RC), precision control (PC) when available, result bits and status
word. Predictions for new hardware tests are fixed before capture. Replays
compare the specified calculation with saved observations; their counts are
reported separately from new instruction executions.

The Xeon and Core i7-6700 test environments report CPU signatures \hex{00050654} and
\hex{000506e3}, with microcode values \hex{0x1} and \hex{0xf0}, respectively.
The Xeon runs under a hypervisor. These are the identifiers reported by the
test environments; we have not independently verified their physical microcode
images. Comparisons between the CPUs use only inputs tested on both.

Separate integer and rational implementations evaluate the same operations
with independent input decoders and quantizers. Exact intermediate values
connect the measured last bits to rounding decisions inside the calculation.
Bounds on reachable values then establish the widths needed by the integer
implementation.

\subsection{How the approximations work}
\label{sec:language}

\subsubsection{Why reduction, polynomials and tables work}

The basic idea is to make the difficult part of the calculation small.
An identity reduces the input to a short interval, where a short polynomial
can approximate the function accurately. Table values, powers of two or
quadrant identities then give the answer for the original input. Harrison,
Kubaska, Story and Tang describe this method and its tradeoffs
\cite[pp.~1--2]{harrison1999}: smaller intervals allow shorter polynomials,
but require more table storage or more work to combine the pieces.

Taylor series explain the powers and leading coefficients. However, matching
a function's derivatives at one point need not give the smallest error over
an interval. For a chosen set $\mathcal P$ of polynomials, one way to choose
the coefficients is
\begin{equation}
 \min_{p\in\mathcal P}\ \max_{r\in I}|f(r)-p(r)|.
\end{equation}
This minimizes the worst error on the interval $I$, a criterion called minimax.
The set $\mathcal P$ can fix the degree, symmetry or leading terms. Harrison
et al. describe how the Remez algorithm finds such polynomials. Shirriff uses
this idea to explain why the Pentium coefficients differ from the reciprocal
factorials and odd reciprocals in Taylor series~\cite{harrison1999,shirriff2025rom}.
The series below explain the form of each approximation. To reproduce the
processor, the code uses the saved ROM coefficients. We have not proved that
every rounded ROM coefficient belongs to an exact minimax solution.

There are three sources of error to keep separate: reducing the input,
approximating the function with a polynomial, and rounding the operations
used to evaluate that polynomial. Harrison treats these separately in his
trigonometric verification~\cite[Sections~4--5]{harrison2000}. Two ways of
evaluating the same polynomial can give different last bits because they
round at different points. The identities explain why each approximation
works; the following rounding rules specify the calculation we reproduce.

\subsubsection{Numbers and rounding}

A normal raw80 value consists of a sign $s$, biased exponent $E$ and explicit
64-bit integer significand $S$:
\begin{equation}
x=(-1)^s S\,2^{E-16383-63}.
\end{equation}
Exponent zero uses effective exponent one for decoding subnormal and
pseudo-denormal values. Each implementation checks the special and unsupported encodings it handles
before starting the ordinary calculation. Supported inputs and status flags
are documented for each instruction.

Let $T_p(v)$ truncate magnitude to $p$ significant bits, preserving sign,
and let $\N{p}(v)$ round to nearest with ties to even. These internal
operators have an unbounded exponent in the reference specification.
Final architectural rounding is denoted $\R(v)$ and uses the requested RC.
Unmarked operations are exact. In particular, a sequence containing
$T_{67}$ followed by $\N{64}$ is not simplified into a single $\N{64}$.

\begin{table}[ht]
\centering\small
\begin{tabular}{@{}>{\raggedright\arraybackslash}p{0.24\linewidth}>{\raggedright\arraybackslash}p{0.42\linewidth}>{\raggedright\arraybackslash}p{0.25\linewidth}@{}}
\toprule
Algorithm & Usual multiplication & Where the rules differ\\
\midrule
Standalone polynomial & $T_{67}(T_{67}(a)T_{64}(b))$ & Ordered input cuts\\
Paired polynomial & $T_{67}(ab)$ & Cut every Horner product before adding\\
Shared table & $T_{67}(T_{67}(a)T_{64}(b))$ & One product rounds directly to RN64\\
\bottomrule
\end{tabular}
\caption{Multiplication rules in the sine and cosine algorithms. The full pseudocode
identifies every operation that uses a different rounding rule.}
\label{tab:operators}
\end{table}

For ordinary results, C1 records whether final rounding increases the magnitude:
\begin{equation}
 C1=\mathbf{1}\{|\R(v)|>|v|\}.
\end{equation}
Special cases and inputs returned without calculation use the rules in the
pseudocode. This C1 rule does not cover every possible initial x87 state or
undefined condition bit. The numerical interface does not predict every
exception flag.

\subsection{Reconstructed algorithms}
\label{sec:trig}

The reconstruction has to determine more than the polynomial coefficients.
It must determine which values enter each operation, when their low bits are
discarded, and how the final correction is combined with the leading term.
We describe the mathematical approximation first, then the finite arithmetic
that matches the tested processors. The complete program is in
Appendix~\ref{app:trig}.


\subsubsection{Using symmetry and nearby table values}

In real arithmetic, writing $x=k\pi/2+r_*$ with $|r_*|\leq\pi/4$ reduces
sine and cosine to a small interval. The quadrant $k\bmod4$ determines the
signs and whether sine and cosine exchange roles. On that interval, sine is odd and cosine is even, giving the forms
\begin{align}
 \sin r&=r+r^3\left(-\frac1{3!}+\frac{r^2}{5!}
                                  -\frac{r^4}{7!}+\cdots\right),\\
 \cos r&=1+r^2\left(-\frac1{2!}+\frac{r^2}{4!}
                                  -\frac{r^4}{6!}+\cdots\right).
\end{align}
Thus the corrections are polynomials in $r^2$, with leading sizes $r^3$
and $r^2$ respectively~\cite[Section~4.19]{dlmfElementary}. Small residuals
make higher powers small: for example, truncating the sine series after
$r^{2n+1}$ gives, for real $|r|\leq1$,
\begin{equation}
 \left|\sin r-\sum_{j=0}^{n}\frac{(-1)^j r^{2j+1}}{(2j+1)!}\right|
 \leq\frac{|r|^{2n+3}}{(2n+3)!}.
\end{equation}
This alternating-series bound illustrates the benefit of reduction; it is
not an error bound for the different ROM coefficients. The same leading
terms explain why sufficiently tiny inputs need only a leading value and
a directed-rounding correction.

For larger residuals, choose a nearby table center $\theta$ and put
$a=r-\theta$. Let $s_\theta=\sin\theta$, $c_\theta=\cos\theta$,
$\sigma=\sin a$ and $\kappa=\cos a-1$. The addition identities become
\begin{align}
 \sin(\theta+a)&=s_\theta+(s_\theta\kappa+c_\theta\sigma),\\
 \cos(\theta+a)&=c_\theta+(c_\theta\kappa-s_\theta\sigma).
\end{align}
Only the small corrections $\sigma$ and $\kappa$ need polynomial evaluation
\cite[Section~4.21]{dlmfElementary}. Harrison et al. use the same addition
identities with different tables~\cite[p.~3]{harrison1999}; the dyadic centers
used here come from the decoded Pentium table~\cite{shirriff2025rom}.
The code uses stored ROM values for $s_\theta,c_\theta$ and rounds the
correction steps as described below.

\subsubsection{Reducing the input angle}

For a finite supported argument, magnitudes at least $2^{63}$ set C2 and
leave the operand unchanged. Below the direct boundary
\hex{3ffe:c90fdaa22168c234}, reduction is bypassed. Otherwise write
$|x|=S2^{e-63}$ and define
\begin{equation}
 D=\hex{3243f6a8885a308d3},\qquad A=S2^{e+2}.
\end{equation}
The integer $D$ represents the fixed divisor $D2^{-65}$. With
$A=qD+t$, $0\leq t<D$, choose $n=q+\mathbf1\{2t>D\}$ and
\begin{equation}
 r=(A-nD)2^{-65}.
\end{equation}
Apply the input sign to both $r$ and $n$. Since $D$ is odd, there is no
halfway integer remainder. Integer division keeps the low bits needed to reduce even a large input
accurately with this divisor.

Exactness here concerns the finite divisor $P=D2^{-65}$, which approximates
$\pi/2$. For the same selected integer $n$, the computed residual $r=x-nP$
and the mathematical residual $r_*=x-n\pi/2$ differ by
\begin{equation}
r-r_*=n(\pi/2-P).
\end{equation}
The phase error can therefore grow with argument magnitude even when the
integer reduction has no arithmetic error. This is the range-reduction
limitation discussed in Intel's comparison with mathematical trigonometric
functions~\cite{intelDifference}; increasing polynomial accuracy cannot
remove it.

The finite reduction has an additional useful property. A reduced nonzero
input cannot have an exactly zero residual: $A=nD$ would imply $D\mid S$, because $D$
is odd and coprime to $2^{e+2}$. But $D$ is larger than a 64-bit significand.
Every reduced residual therefore lies on the nonzero $2^{-65}$ lattice.
This gives a concrete domain for the subsequent tiny and polynomial paths,
and lets tests work backward from a desired residual to possible raw80 inputs.

For sine and cosine, the residual's magnitude chooses the tiny-input method,
polynomial or table, with boundaries at $2^{-32}$ and $1/4$. Table centers have numerators
$18,22,26,30,36,44,52$ over 64. Selection and center subtraction remain
exact at the cell boundaries. The quadrant determines which internal function supplies each external
result. If the requested phase is $\phi=0$ for sine or $\phi=1$ for cosine,
use $(n+\phi)\bmod4$ to choose the internal sine/cosine lane and its sign.
Apply that sign before final rounding; rounding a positive magnitude and
negating it afterwards gives the wrong directed result.

\subsubsection{Standalone polynomial: operand widths and two coefficient chains}
\label{sec:standalone}

The standalone polynomial separates odd and even coefficient indices and
uses the asymmetric ordinary product in Table~\ref{tab:operators}.
Let $K_1,\ldots,K_6$ be the appropriate sine or cosine coefficients,
$M(a,b)=T_{67}(T_{67}(a)T_{64}(b))$, and $A(a,b)=\N{64}(a+b)$.
The program forms $u=M(r,r)$ and $v=M(u,u)$, then
\begin{align}
 o&=A(K_1,M(v,A(K_3,M(v,K_5)))),\\
 e&=A(K_2,M(v,A(K_4,M(v,K_6)))),\\
 \ell&=M(u,o),\qquad h=M(v,e).
\end{align}
Before final rounding, sine is $r+M(r,A(\ell,h))$ and cosine is
$1+T_{67}(\ell+h)$. The different final operations affect the result bits.

The asymmetric product is particularly informative at the fourth power.
The first square can contain 67 significant bits. Forming $v=M(u,u)$ then
uses the full retained square on one side and its 64-bit truncation on the
other. Replacing that operation with a symmetric $T_{67}(u^2)$ can change
the polynomial even though both descriptions appear to use ``67-bit
multiplication.''

The reachable widths explain where this distinction matters. A direct
polynomial residual has at most 64 significant bits; a reduced polynomial
residual has at most 63, because it is below $1/4$ on the $2^{-65}$ lattice.
The coefficient widths and RN64 additions make the other input cuts
identities under the displayed operand order. The fourth-power operation is
the only nonidentity input cut in this polynomial graph. This is a statement
about the numerical program; it does not establish physical multiplier ports.

The two terminal expressions matter just as much. Sine rounds $\ell+h$ to
RN64, multiplies by $r$, and adds that correction to $r$. Cosine truncates
$\ell+h$ to 67 bits before adding it to one. Sharing the coefficient-chain
code does not make those endings interchangeable. Internal rounding remains
fixed when the caller changes RC; the requested mode acts on the signed final
value.

\subsubsection{FSINCOS: a different evaluation order}
\label{sec:paired}

Paired \ins{FSINCOS} instead shares $u=T_{67}(r^2)$ between two Horner
chains. At every coefficient stage, each chain uses
\begin{equation}
 p_i=\N{64}\!\left(T_{67}(p_{i+1}u)+S_i\right),
 \quad
 q_i=\N{64}\!\left(T_{67}(q_{i+1}u)+C_i\right).
\end{equation}
Before final rounding, paired sine is $r+T_{67}(\N{64}(p_1u)r)$ and cosine
is $1+T_{67}(q_1u)$. Calling FSIN and FCOS separately therefore does not
reproduce FSINCOS. Its final C1 describes rounding of the returned cosine,
even when the quadrant makes the internal sine calculation supply that result.

The Horner stages operate on the small correction polynomials in $r^2$.
Truncation to 67 bits retains three more significant bits than each RN64
coefficient sum. Those bits participate in the next addition before its
rounding. Measurements near final rounding midpoints constrain this order;
exact arithmetic traces the low bits through each stage. The same product
cut applies at all five coefficient additions on both arms.

\subsubsection{Shared table reconstruction: where a product is rounded}
\label{sec:trig-table}

For $r\geq1/4$, select a center $b/64$ by
\begin{equation}
 b=\begin{cases}
 18+4\lfloor16(r-1/4)\rfloor,&r<1/2,\\
 36+8\min(2,\lfloor8(r-1/2)\rfloor),&r\geq1/2.
 \end{cases}
\end{equation}
The reachable numerators are $18,22,26,30,36,44,52$. Set $a=r-b/64$
and $u=M(a,a)$. With the standalone operators $M$ and $A$, define the
four-coefficient chain
\begin{equation}
 H(K)=A(K_1,M(u,A(K_2,M(u,A(K_3,M(u,K_4)))))).
\end{equation}
Let $p=H(S_4)$ and $q=H(C_4)$. The highest-degree sine coefficient includes
the established adjustment $-2^{-25}$ to the decoded constant. This adjustment
is explicit in the specification; the other coefficients retain their ROM
values.

The small sine and cosine corrections are
\begin{equation}
 v=A(a,M(M(u,p),a)),\qquad
 w=M_R(u,q),\qquad M_R(x,y)=\N{64}(T_{67}(x)T_{64}(y)).
\end{equation}
Using the stored center values $t_s,t_c$, form
\begin{align}
 z_s&=T_{67}\bigl(M(t_c,v)+M(t_s,w)\bigr),\\
 z_c&=T_{67}\bigl(-M(t_s,v)+M(t_c,w)\bigr).
\end{align}
The positive-residual prevalues are $t_s+z_s$ and $t_c+z_c$. Quadrant sign
restoration precedes the final RC rounding. These are the addition identities
above with explicit rounding at each arithmetic step.

The definition of $M_R$ rounds the exact product of its cut operands directly
to RN64. Its rounding decision uses the product's low bits, including those
below the 67-bit position. By contrast, ordinary $M$ products discard those
bits. The destination of each product is therefore part of the algorithm.

Exact cell selection matters independently of polynomial rounding. Converting
the input to binary64 before classifying it merges raw80 values immediately
below a boundary with the boundary itself. The classifier and subtraction
therefore operate on the original dyadic value. Center subtraction leaves
at most 61 significant offset bits for reduced inputs and 60 for direct
inputs. All table input cuts are identities under the stated operand order;
the table tests establish the rounding destinations, without separately
establishing an asymmetric physical input port.

At an exact center the correction vanishes, exposing the ROM entry and final
rounding directly. Center measurements therefore constrain the stored table values
independently of the correction polynomial. The center $52/64$ itself lies
outside the reachable residual interval, although its cell is used. The ROM's $60/64$ row is never selected by this
dispatcher.

\subsubsection{Tiny results}

The external-input bypass at $|x|<2^{-68}$ returns the input for sine
and $+1$ for cosine, with its specified C1 convention. In the remaining
tiny-residual region, the leading magnitude is $r$ or $1$. RN and rounding
away from zero retain that leading magnitude; rounding toward zero selects
its predecessor at 64 significant bits. When the leading value is a power of two, that predecessor has an exponent
one less than the leading value. This path and
the external-input bypass must remain distinct.

For polynomial and table results, C1 reports a final magnitude increment.
For tiny non-bypass results, it is one when the leading magnitude is retained
and zero when the predecessor is selected. The external-input bypass returns
C1=0 in every mode. These rules specify both the returned value and the status of its final
rounding.

\subsubsection{From exact arithmetic to the integer implementation}

The executable specification uses signed dyadic values: an integer multiplied
by a power of two. Multiplication and aligned addition are exact until an
explicit quantizer discards bits. The C implementation uses bounded integer
carriers whose widths follow from bounds on the reachable values.

The reducer's integer division and nonzero-residual argument establish its
input domain. Polynomial and table width bounds then cover their products,
sums and final rounding. For the paired polynomial, exact interval enclosures
cover all 30 polynomial binades. Its 840 operation-bound checks require at
most 139 accumulator bits and alignment shifts at most 132, within the
signed 256-bit carrier. Separate conversion and dispatch checks establish
that supported external inputs reach the intended calculation.

These arguments support the specified implementation under its entry and
quantizer contracts. They complement the independently written rational
program. Hardware measurements provide the separate evidence connecting
that calculation to the processor.

\subsection{Sine and cosine constants}
\label{sec:trig-constants}

The decoded ROM supplies the six- and four-coefficient polynomial families
and the sine/cosine table. Appendix~\ref{app:trig-constants} lists their
literal values and ROM row numbers, with the two differences from the saved
cosine-table transcription. The fixed $-2^{-25}$ sine-coefficient adjustment
appears explicitly in the algorithm.

\subsection{Testing the sine and cosine reconstruction}
\label{sec:validation}

\subsubsection{How measurements constrain the calculation}

The reconstruction separates the calculation into reduction, polynomial or
table evaluation, and final rounding. Each offers a way to expose otherwise
hidden arithmetic. Exact table centers isolate the stored values. Inputs on
either side of an interval boundary identify dispatch rules. Values near
output rounding boundaries constrain the low bits retained by intermediate
operations. Raw80 inputs supply significand bits that a binary64 generator
cannot exercise.

For example, let $y$ be an ordinary observed result under RN, with adjacent
representable values $y_-$ and $y_+$. The signed value $v$ just before final
rounding must lie in its rounding interval:
\begin{equation}
 \frac{y_-+y}{2}\ \leq\ v\ \leq\ \frac{y+y_+}{2},
\end{equation}
with midpoint inclusion determined by ties to even. C1 adds a magnitude
constraint: C1=1 requires $|v|<|y|$, while C1=0 requires $|v|\geq|y|$.
The directed modes provide further interval bounds. Exact evaluation of the
polynomial and table operations connects these observable constraints to
intermediate widths, operation order and rounding destinations.

The reducer also lets us choose an internal residual and find external
inputs that reach it. For a desired residual integer $d$, solve
$S2^{e+2}-nD=d$ subject to the significand and exponent bounds. An exact
preimage repeats that residual at a different quadrant or magnitude; nearby
external inputs bracket it. Sign counterparts and all four RC settings
then check sign restoration and final rounding while preserving the kernel.

The structural tests cover direct and reduced polynomial inputs, table cells
and reachable centers, tiny-input boundaries and encoding classes. The
exact-center set covers 672 specified combinations: 28 signed external
operands, two instructions, four RC settings and three PC settings. Broad
samples and archive replays provide complementary coverage beyond these
selected boundaries.

\subsubsection{Completed native runs and archive checks}

Table~\ref{tab:trig-evidence} includes the large completed trig tests, later
boundary tests and independent reference replays. Its counts need some care:
an instruction execution, a saved row checked again in software and an
output value are different things. FSINCOS produces two output values.
Saved test sets overlap, and checking the same data on another host does not
add new hardware observations.

\input{generated-suite/trig-evidence-table.tex}

The completed native runs execute FSIN \TrigNativeRows{} times on the
Xeon. They test \TrigNativeInputs{} inputs drawn from binary64 encodings,
using round-to-nearest, downward and upward rounding at PC64. All reported
result, C1 and interval-check mismatch
counters are zero. C1 is checked on non-C2 returns using the model's
upward and downward rounded results. The sample includes special values
and out-of-range returns. These are selected binary64 encodings, not an
exhaustive input set. The retained evidence consists of completion logs
and a build record. Individual outputs and a separate hash of the test
executable are unavailable for these runs.

The standalone and paired archive replays check \TrigIsevenRows{} saved
rows on the i7 host and \TrigXeonRows{} on the Xeon host. They contain
356,459,555 and 29,648,279 output-value comparisons, respectively, and
193,825,064 and 17,600,228 applicable recorded C1 checks. All checked
results and flags agree. Each host checks a common 15,655,784-row set
plus its available archive files. The replay host does not necessarily
identify the CPU that captured a row; unrecorded modes and flags are
excluded. Overlapping test sets are not added to these counts.

A separate hardware challenge tests 6,180 operands on each CPU with three
instructions, four RC settings and three PC settings: \TrigChallengeRows{}
executions and 296,640 output-value comparisons per CPU. The implementation
and predictions were fixed before capture. The evidence register identifies
the source versions, input sets, recorded fields and CPU contexts for the
native runs, archive replays and this challenge.

\subsubsection{What the checks leave open}

Independent rational arithmetic checks the polynomial and table stages with
separate decoders and quantizers; a separate integer formula checks the tiny
path. These checks support the implementation of the specified operations.
The hardware agreement covers the reported inputs and CPU contexts. It does
not establish agreement for every raw80 input or uniquely identify the
processor's physical implementation.

The numerical entry points do not model every incoming FPU state or unmasked
trap. Undefined status fields and unrecorded flags receive no pass credit.
Testing a Core i7 and a Xeon establishes the reported agreement on those
Skylake contexts; it does not establish identical behavior across all Intel
generations. The separate Itanium implementation remains a historical
comparison, not an interchangeable source of expected Skylake answers.

\subsection{Result and implementation}

The main result is a fixed arithmetic description of FSIN, FCOS and FSINCOS
that explains their recorded last-bit behavior. Public ROM constants give
the approximation data. The reconstruction supplies the standalone and
paired evaluation orders, operand-width rules, table rounding destinations,
and final sign/C1 behavior needed to turn that data into matching results.
Algebraic equivalence is insufficient at each of those points.

The integer implementation exposes separate raw80 entry points for standalone
and paired instructions. The programmer guide describes the command-line
interface, supported inputs and replay tools.

\clearpage
\section{FPTAN and FPATAN}
\label{family:tangent}

Tangent and arctangent form a separate, shorter part of the reconstruction.
Their polynomial and table methods followed comparatively directly from
Shirriff's published ROM account~\cite{shirriff2025rom}. This section gives
their algorithms and evidence together; Appendix~\ref{app:tangent} contains
the complete executable references.

\subsection{FPTAN: retain the internal pair before division}

The identity $\tan x=\sin x/\cos x$ explains the tangent method. FPTAN
uses the finite angle divisor specified in Section~\ref{sec:trig}, selects a direct polynomial or table
correction, and retains an internal sine/cosine pair. Ordinary products
truncate to 67 bits and coefficient additions round to RN64; specified
scaling products use RN64. The internal pair is cut to 67 bits, transformed
by the quadrant, and divided before final RC rounding. Dividing architectural
FSIN and FCOS results would introduce different rounding steps.

Near a pole, a small denominator makes the quotient sensitive to errors in
that pair. The same finite-divisor limitation as sine/cosine also applies.
The direct tiny bypass returns the operand. Ordinary successful calls push
$+1$; range returns leave the operand and stack depth unchanged. The reference
specifies the special-result pushes and its result/C1/C2 contract.

\subsection{FPATAN: reduce a ratio to a small angle}
\label{sec:atan}

After ordering the nonzero magnitudes as $a\leq b$, let $r=a/b$. The direct
path uses a small-ratio polynomial. The table path uses
\begin{equation}
 \arctan r=\arctan c+\arctan t,\qquad
 t=\frac{r-c}{1+rc}=\frac{a-cb}{b+ca}.
\end{equation}
Centers spaced by $1/32$ give $|t|\leq1/64$. The series
$\arctan t=t-t^3/3+t^5/5-\cdots$ explains the correction polynomial
\cite[Section~4.24]{dlmfElementary}. The code forms the residual from $a,b$
with the specified 67-bit cuts, evaluates interleaved coefficient chains,
and adds the stored angle. Magnitude ordering and operand signs determine
the final $\pi/2$ or $\pi$ adjustment before architectural rounding.
The wrapper also handles signed zeros, infinities, NaNs and masked exceptions.
Its direct square uses $\N{64}(zT_{64}(z))$. The full reference specifies
every product width and rounding destination.

\subsection{Constants and published implementation evidence}

FPTAN uses the sine/cosine coefficient and table values already listed in
Appendix~\ref{app:trig-constants}, including the same two cosine-table transcription differences. FPATAN has its own ratio-polynomial coefficients, angle table and
quadrant constants, listed in Appendix~\ref{app:atan-constants}.

The arctangent audit also finds 42 matching constant bit patterns and the
coefficient-read order in public Goldmont ROM material
\cite{ucodeDisasm,customProcessingUnit}. Those comparisons cover the exposed
bits; the remaining low bits, signs, exponents and operation meanings need
separate analysis. They guide the tested algorithm without establishing a
recovered Skylake microcode listing.

\subsection{Validation and scope}
\label{sec:tangent-validation}

\input{generated-suite/tangent-evidence-table.tex}

The FPTAN reference checks both complete T0002/T0003 captures, with
\TanSuccess{} successful calls and \TanRange{} C2 returns per CPU. It checks
the tangent, pushed value, C1/C2, controls and stack transitions. The final
set contains 122,900 normal finite raw80 operands and \TanRows{} input/control
tuples per CPU. It includes rounding-boundary, modular-residue, exponent and
neighboring-input tests. Nine neighboring-input windows have complete coverage. Captured status words agree between CPUs, but the model
does not predict every exception flag.

The FPATAN reference matches \AtanRows{} saved tuples in \AtanPacks{} packs,
checking results, C1, arithmetic exceptions and operand-preload flags.
All catalog tuples have matching captured results and status on both CPUs. The catalog includes the 445,588-observation independent challenge per
CPU, with mathematical-boundary and full-significand input selection.
Exact table-index equivalence and tests of RN64 ties are documented in its
separate algorithm and evidence records.

The FPTAN CLI returns raw80 values and range results. The FPATAN CLI and
GMP-based C API take ordered $y,x$ inputs with explicit RC/PC and return a
raw80 angle, C1 and modeled masked exceptions. The wrapper assumes clear
initial exception latches and a valid two-deep stack. The captured status
agreement and the model's predicted fields remain separate claims.

\clearpage
\section{F2XM1, FYL2X and FYL2XP1}
\label{family:explog}

The exponential and logarithm implementations complete the reference.
Their approximation families follow directly from the published ROM material
\cite{shirriff2025rom}.
The short descriptions below give the mathematical identities and the
specified rounding. Appendix~\ref{app:explog} contains the full programs.

\subsection{F2XM1: compute the small exponential difference directly}
\label{sec:exp}

For $t=x\ln2$, the series
\begin{equation}
 2^x-1=t+t^2\left(\frac12+\frac{t}{6}+\frac{t^2}{24}+\cdots\right)
\end{equation}
explains the linear and polynomial paths~\cite[Section~4.2]{dlmfElementary}.
Computing the difference directly preserves a small result that subtracting
one from an already rounded exponential could erase. The tiny path applies
below $2^{-68}$; an eleven-coefficient polynomial applies up to $|x|=1/4$.
Above that boundary, a stored $2^b-1$ and a shorter correction use
\begin{equation}
 2^x-1=(2^b-1)+2^b\bigl(\exp((x-b)\ln2)-1\bigr).
\end{equation}
The selected centers bound $|x-b|$ by $1/64$. The exact reference preserves
the distinct RN64 and 67-bit versions of the scaling products. The endpoints
$x=1,-1$ return $1,-1/2$. Final rounding, including a subnormal tiny result,
uses raw80 spacing and the requested RC.

\subsection{FYL2X and FYL2XP1: a logarithm followed by multiplication}
\label{sec:log}

Both instructions retain an internal logarithm $g$, then round $yg$ to
raw80. Their direct and table paths follow
\begin{equation}
 \ln v=2\operatorname{atanh}t
       =2\left(t+\frac{t^3}{3}+\frac{t^5}{5}+\cdots\right),
 \qquad t=\frac{v-1}{v+1}.
\end{equation}
Reducing $v$ toward one makes the odd-power correction small
\cite[Section~4.6]{dlmfElementary}. The direct intervals are
$7/8\leq x\leq9/8$ for FYL2X and $|x|\leq1/8$ for FYL2XP1.
For FYL2XP1 the transformed argument is $x/(2+x)$, so the calculation
preserves the increment without first rounding $1+x$ to 64 bits.

Otherwise, normalization $x=m2^e$ and a midpoint $a_i=(65+2i)/64$ give
\begin{equation}
 \log_2x=e+\log_2a_i+
 2\log_2(\mathrm e)\operatorname{atanh}\!\left(\frac{m-a_i}{m+a_i}\right).
\end{equation}
The 32 mantissa cells bound the exact transformed argument by $1/129$.
Each table logarithm is split into a 40-bit high part and a 67-bit low
correction. The direct and table polynomials, retained additions and final
product use the instruction-specific rounding given in the reference.
FYL2XP1 outside its direct interval enters the table program with $T_{67}(1+x)$.
Its supported domain is Intel's guaranteed interval. The numerical wrappers
include C1 and masked exceptions under their documented initial-state contract.

\subsection{Constants and source checks}
\label{sec:explog-constants}

Appendix~\ref{app:explog-constants} lists the exponential and logarithm
constants and the five literal differences from the saved published
transcription. The logarithm high/low splits derived at decimal precisions
128 and 192 agree on every word. The upper 64 bits of its 77 constants
match the saved public Goldmont ROM, as do the two polynomial blocks and
split-table read order~\cite{ucodeDisasm,customProcessingUnit}. These
comparisons cover the data and operations exposed by those sources.

\subsection{Validation and interfaces}
\label{sec:explog-validation}

\input{generated-suite/explog-evidence-table.tex}

The F2XM1 reference matches \FtwoRows{} saved results and C1 values.
The raw80 implementation also matches \FtwoRawRows{} cases on each
processor, covering selected subnormal, pseudo-denormal and minimum-normal
families across all RC settings and selected complete PC groups. These
comparisons check numerical results and C1.

The logarithm C and rational models match \LogRows{} saved observations,
including \LogXRows{} FYL2X and \LogPoneRows{} FYL2XP1 cases. The
\LogFreshRows{}-case hardware pack used the C source fixed before capture.
A subsequent i7 confirmation matches 51,636 selected Xeon tuples; the other
890,172 Xeon tuples are outside that confirmation. The checks distinguish
instruction-specific PE/UE rules and operand-preload effects from final
result rounding.

F2XM1 uses the unary raw80 CLI. The logarithms also provide GMP-based C
APIs taking ordered $y,x$ inputs and explicit RC/PC. Their result includes
C1 and modeled masked exceptions under the documented initial-state contract.
The offline package includes 216 F2XM1 boundary regressions and 854 logarithm
C/API/reference examples. These checks preserve the implemented boundary
rules without requiring native x87 execution.

\paragraph{Availability and acknowledgments.}
The accompanying package contains the paper source, programmer documentation,
complete canonical references and verification records. Its build and small
saved-result checks run without native x87; repeating the large archive
replays requires the separately saved data. File hashes identify the exact
source, constants, evidence and reviewed PDF. Public repository and
archival identifiers have not been assigned in this version. Licensing of
the author's original work remains undecided; third-party source notices
are preserved and attributed separately. This work builds directly on Ken
Shirriff's published ROM research and the cited public implementations and
microcode tools. Substantial AI assistance, including Anthropic Claude and
OpenAI Codex, was used in the investigation and drafting. The author is
responsible for the measurements, code and claims.

\bibliographystyle{plain}
\begingroup\small\raggedright
\bibliography{suite-references}
\endgroup

\clearpage
\appendix
\section{FSIN, FCOS and FSINCOS reference}
\label{app:trig}

\subsection{Complete pseudocode}
The pseudocode below specifies the complete numerical calculation.
Arithmetic is exact except at named rounding operations. The arrays use
one-based coefficient indices; the literal inventory below retains its
original zero-based labels. Complete arbitrary x87 state is outside this
numerical walkthrough.
\input{generated-suite/trig-listings.tex}

\subsection{Sine and cosine constants}
\label{app:trig-constants}
Each row represents $(-1)^s S2^q$, with hexadecimal integer $S$ and its
original ROM row number. The sine-coefficient adjustment is explicit in
the algorithm. The two cosine-table differences from the saved published
transcription are recorded below, followed by the complete literal inventory.

\begin{table}[ht]
\centering\footnotesize
\begin{tabular}{@{}rlll@{}}
\toprule
ROM row & Role & Saved published significand & Used significand\\
\midrule
186 & $\cos(44/64)$ & \hex{62ec41e9772401864} & \hex{62ec4169772401864}\\
189 & $\cos(18/64)$ & \hex{7af8853ddbbe9ffd0} & \hex{7af8853ddbbe9efd0}\\
\bottomrule
\end{tabular}
\caption{The two corrected cosine-table words used by the sine/cosine reconstruction.}
\end{table}

\input{generated-suite/trig-literals.tex}

\clearpage
\section{FPTAN and FPATAN reference}
\label{app:tangent}

\subsection{FPTAN executable calculation}
These functions from \path{docs/sibling_reference.py} use the arithmetic and
dispatch helpers in Appendix~\ref{app:common}. Constant arrays are
zero-based and use the sine/cosine inventory in
Appendix~\ref{app:trig-constants}. The demonstrated domain and predicted
fields are recorded in Section~\ref{sec:tangent-validation}.
\input{generated-suite/tangent-listings.tex}

\subsection{FPATAN executable reference}
\label{app:atan}
The following blocks specify the FPATAN numerical calculation and wrapper. \ins{ROM} maps the literal row numbers below to exact rational
values. The wrapper assumes masked exceptions and clear initial latches
with a valid two-deep stack.
\input{generated-suite/atan-listings.tex}

\subsection{Arctangent constants}
\label{app:atan-constants}
Rows retain their original ROM identities and represent $(-1)^s S2^q$.
The tangent calculation uses the previously listed sine/cosine constants.
\input{generated-suite/atan-literals.tex}

\clearpage
\section{F2XM1, FYL2X and FYL2XP1 reference}
\label{app:explog}

\subsection{F2XM1 executable kernels}
These functions from \path{docs/sibling_reference.py} use the shared
decoder, arithmetic, raw80 rounder and entry dispatch in
Appendix~\ref{app:common}. The tiny path rounds directly to raw80.
\input{generated-suite/exponential-listings.tex}

\subsection{Logarithm executable reference}
\label{app:log}
The independent rational program below is copied from
\path{fyl2x-re/model.py}. Its default \ins{Policy} specifies the numerical
schedule implemented by the C CLI and API. The ROM is read with the four
logarithm-table transcription differences listed below.
\input{generated-suite/log-listings.tex}

\subsection{Exponential and logarithm constants}
\label{app:explog-constants}
Each literal represents $(-1)^s S2^q$ and retains its original ROM row.
The five differences from the saved published transcription are recorded
below. The three family inventories contain 239 entries in total.

\begin{table}[ht]
\centering\footnotesize
\begin{tabular}{@{}rlll@{}}
\toprule
ROM row & Role & Saved published significand & Used significand\\
\midrule
110 & $2^{-57/128}-1$ & \hex{43fcb5810d1604f37} & \hex{43fcb5810d1604f33}\\
207 & log table high, $i=1$ & \hex{43ace37e8a8000000} & \hex{43ace27e8a8000000}\\
245 & log table low, $i=7$ & \hex{6f66cc899b64303f7} & \hex{6f66cc899b64b03f7}\\
259 & log table low, $i=21$ & \hex{5c833d56efe4338fe} & \hex{5c813d56efe4338fe}\\
268 & log table low, $i=30$ & \hex{4a83eb6e0f93f7a64} & \hex{4a83eb6e0f93f7a44}\\
\bottomrule
\end{tabular}
\caption{Corrected exponential and logarithm table words in the saved published transcription.}
\end{table}

{\renewcommand{\arraystretch}{0.94}
\input{generated-suite/explog-literals.tex}
}

\clearpage
\section{Shared executable-reference support}
\label{app:common}

The FPTAN and F2XM1 references share a Python source file. The functions
below supply its exact arithmetic, decoding, storage and entry dispatch.
The complete file is included in the programmer package.
Unsupported raw encodings are outside this reference's contract.
Constants are loaded as exact dyadics from
\path{docs/sibling-constants.json}.
\input{generated-suite/sibling-common-listings.tex}

\end{document}
